\documentclass[12pt]{article}
\usepackage[utf8]{inputenc}
\usepackage[T1]{fontenc}
\usepackage{hyperref}
\title{Teke, Ibali Language Resources for K-12 Curriculum Development}
\author{Compiled by Claude Sonnet 5 (claude-sonnet-5) for Andr\'as Kornai}
\date{}
\begin{document}
\maketitle
\section{Introduction}
Ibali Teke (ISO 639-3 \texttt{tek}; also Teke-Ibali, Eastern Teke, Kiteke) is a Bantu language spoken in the Democratic Republic of the Congo (Maluku district of Kinshasa; Kwamouth territory, Mai-Ndombe) and in the Republic of the Congo. Speaker counts in the sources found vary widely (roughly 71,000 to 222,000), well above the figure in the task description. Very little educational material was found in search results.
\section{Primary Educational Resources}
No primers, textbooks or school programmes in Ibali Teke were found in search results.
\section{Linguistic and Reference Materials}
\begin{itemize}
\item Wikitongues language page (identification: ISO code, Glottocode, script, vitality): \url{https://wikitongues.org/languages/tek/}
\item Glottolog language entry (classification): \url{https://glottolog.org/resource/languoid/id/bali1275.bigmap.html}
\item Wikipedia overview of East Teke (dialects, speaker estimates): \url{https://wikipedia.com/wiki/ISO_639:tek}
\end{itemize}
\section{Religious and Community Resources}
\begin{itemize}
\item Joshua Project, Eastern Teke in DR Congo (reports Bible portions from 1889--1905, no complete Bible or New Testament): \url{https://joshuaproject.net/index.php/people_groups/11707/CG}
\item Joshua Project, language page: \url{https://www.joshuaproject.net/index.php/languages/tek}
\end{itemize}
\section{Digital and Audio Resources}
\begin{itemize}
\item Global Recordings Network, Kiteke audio (Bible-story series; the page notes no recordings in the Bali variety): \url{https://globalrecordings.net/en/language/tek}
\end{itemize}
\section{Curriculum Development Context}
No evidence of school use was found. Any curriculum would need to start from basic documentation: orthography, a word list, and recorded texts, ideally with community partners in Kinshasa/Maluku and Mai-Ndombe. Related Kiteke materials should be checked against the Ibali variety before use.
\end{document}